\documentclass[../DoAn.tex]{subfiles}
\begin{document}

Chương 3 đã trình bày các công nghệ nền tảng. Chương này trình bày cách thiết kế
và hiện thực hệ thống dựa trên các công nghệ đó, gồm: kiến trúc tổng thể, thiết
kế các mô-đun chính, máy trạng thái điều khiển phương tiện, hai luồng dữ liệu
giữa mô phỏng và hiển thị, cùng kết quả thực nghiệm và đánh giá. Nội dung tập
trung vào những điểm phản ánh đúng hiện thực của hệ thống thay vì mô tả lý thuyết
chung.

\section{Thiết kế kiến trúc}

\subsection{Lựa chọn kiến trúc phần mềm}

Hệ thống được xây dựng theo kiến trúc phân tầng kết hợp mô hình mô phỏng--hiển
thị, trong đó phần mô phỏng giao thông chạy trên \gls{SUMO} (điều khiển bằng
\gls{TraCI} trong Python) và phần hiển thị, tương tác chạy trên Unity. Hai phần
giao tiếp qua kết nối TCP theo cơ chế truyền thông điệp dạng \gls{JSON}, có dấu
hiệu kết thúc gói tin \texttt{<END>} để tách thông điệp trên luồng TCP.

Kiến trúc này được chọn vì các lý do: tách biệt trách nhiệm giữa mô phỏng và hiển
thị; dễ mở rộng thêm loại đối tượng hoặc chỉ số thống kê mà ít ảnh hưởng phần còn
lại; và tương thích với công cụ, do \gls{SUMO} đã có \gls{TraCI} hỗ trợ điều
khiển theo bước, còn Unity thuận lợi cho hiển thị thời gian thực và xử lý vật lý.
Ở mức khái niệm, có thể ánh xạ theo mô hình Model--View--Controller: phần Model là
trạng thái mô phỏng trong \gls{SUMO} cùng các tệp mạng và tuyến; phần View là các
đối tượng và mô-đun dựng cảnh trong Unity; phần Controller là lớp giao tiếp TCP
và bộ xử lý lệnh điều khiển.

\subsection{Kiến trúc tổng quan}

Hệ thống gồm ba tầng: tầng mô phỏng (Python kết hợp \gls{SUMO}/\gls{TraCI}) khởi
tạo mạng, chạy mô phỏng theo bước, trích xuất trạng thái và ghi kết quả; tầng
tích hợp (TCP cùng định dạng \gls{JSON}) đóng gói, tách khung và truyền thông
điệp; tầng hiển thị (Unity) nhận dữ liệu, dựng và cập nhật đối tượng, cung cấp
giao diện điều khiển và xử lý tương tác lái xe.

Điểm đáng chú ý trong hiện thực là hệ thống dùng ba kênh TCP tách theo mục đích
và hai luồng dữ liệu ngược chiều, như minh họa ở Hình \ref{fig:arch-detail}.
Kênh chính truyền dữ liệu mô phỏng từ Python sang Unity; kênh cập nhật nhận dữ
liệu phương tiện do người dùng lái từ Unity về Python; kênh điều khiển nhận các
lệnh tạm dừng, tiếp tục, kết thúc và cập nhật tham số.

\begin{figure}[H]
	\centering
	\begin{tikzpicture}[every node/.style={font=\small}, node distance=1.0cm]
		\node[draw, rounded corners, align=center, minimum height=1cm] (sumo) {SUMO};
		\node[draw, rounded corners, align=center, minimum width=3cm, minimum height=1.6cm, right=2.2cm of sumo] (py) {Python\\(điều phối)};
		\node[draw, rounded corners, align=center, minimum width=3cm, minimum height=1.6cm, right=4.2cm of py] (unity) {Unity\\(hiển thị \& tương tác)};
		% TraCI
		\draw[{Latex}-{Latex}] (sumo) -- node[midway, above, font=\scriptsize] {TraCI} (py);
		% 3 kenh TCP
		\draw[-{Latex}] ([yshift=14pt]py.east) -- ([yshift=14pt]unity.west)
			node[midway, above, font=\scriptsize] {5050: dữ liệu mô phỏng};
		\draw[{Latex}-] ([yshift=-2pt]py.east) -- ([yshift=-2pt]unity.west)
			node[midway, above, font=\scriptsize] {5053: vị trí xe người lái};
		\draw[{Latex}-] ([yshift=-16pt]py.east) -- ([yshift=-16pt]unity.west)
			node[midway, below, font=\scriptsize] {5054: lệnh điều khiển};
	\end{tikzpicture}
	\caption{Kiến trúc chi tiết với ba kênh TCP và hai luồng dữ liệu (bản nháp)}
	\label{fig:arch-detail}
\end{figure}

\subsection{Thiết kế các mô-đun chính}

Phía Python đảm nhiệm mô phỏng và điều phối, gồm các mô-đun chính: mô-đun điều
phối khởi tạo các kênh TCP và chạy vòng lặp mô phỏng; mô-đun mạng đóng gói và
tách khung thông điệp; nhóm mô-đun đọc trạng thái chuyển dữ liệu xe, đèn tín hiệu
và người đi bộ từ \gls{TraCI} sang định dạng chuẩn; mô-đun cập nhật phương tiện
xử lý dữ liệu xe do người dùng lái gửi lên; nhóm mô-đun sinh mạng và tuyến tạo
kịch bản mô phỏng; và mô-đun lưu trữ phiên ghi kết quả của mỗi lần chạy.

Phía Unity đảm nhiệm hiển thị và tương tác, gồm các nhóm mô-đun chính trình bày
trong Bảng \ref{tab:unity-modules}.

\begin{table}[H]
	\centering
	\caption{Các nhóm mô-đun chính phía Unity}
	\label{tab:unity-modules}
	\begin{tabular}{p{0.26\textwidth} p{0.64\textwidth}}
		\hline
		\textbf{Nhóm mô-đun} & \textbf{Nhiệm vụ} \\
		\hline
		Mạng (Network) & Kết nối TCP, nhận dữ liệu mạng đường và dữ liệu mô phỏng. \\
		Giao thông (Traffic) & Quản lý vòng đời và cập nhật phương tiện, đèn tín hiệu, người đi bộ; nội suy chuyển động. \\
		Phương tiện người lái (UnityVehicle) & Chiếm quyền, lái bằng vật lý, gửi dữ liệu xe lên server, xử lý va chạm. \\
		Tối ưu (Optimization) & Lọc đối tượng theo vùng quan sát, giảm số lệnh vẽ. \\
		Giao diện (UI) & Điều khiển camera, tạm dừng, tốc độ, nút chiếm/trả quyền. \\
		Mô phỏng (Simulation) & Điều phối chế độ thời gian thực và chế độ phát lại. \\
		\hline
	\end{tabular}
\end{table}

Cách phân chia mô-đun và quan hệ phụ thuộc giữa hai phía được minh họa ở Hình
\ref{fig:packages}.

\begin{figure}[H]
	\centering
	\begin{tikzpicture}[every node/.style={font=\footnotesize},
		pkg/.style={draw, rounded corners=1pt, minimum width=3.4cm, minimum height=0.7cm,
			align=center, fill=black!2}, >=Latex, node distance=0.4cm and 3.4cm]
		% Python packages
		\node[pkg] (p1) {Điều phối};
		\node[pkg, below=of p1] (p2) {Mạng (tách khung)};
		\node[pkg, below=of p2] (p3) {Đọc trạng thái};
		\node[pkg, below=of p3] (p4) {Cập nhật phương tiện};
		\node[pkg, below=of p4] (p5) {Sinh mạng \& tuyến};
		\node[pkg, below=of p5] (p6) {Lưu phiên};
		% Unity packages
		\node[pkg, right=of p1] (u1) {Network};
		\node[pkg, below=of u1] (u2) {Traffic};
		\node[pkg, below=of u2] (u3) {UnityVehicle};
		\node[pkg, below=of u3] (u4) {Optimization};
		\node[pkg, below=of u4] (u5) {UI};
		\node[pkg, below=of u5] (u6) {Simulation};
		% group boxes
		\node[draw, dashed, rounded corners, fit=(p1)(p6), inner sep=10pt,
			label=above:{\textbf{Python --- Mô phỏng \& điều phối}}] (gp) {};
		\node[draw, dashed, rounded corners, fit=(u1)(u6), inner sep=10pt,
			label=above:{\textbf{Unity --- Hiển thị \& tương tác}}] (gu) {};
		% dependency
		\draw[{Latex}-{Latex}] (gp.east) -- node[above, font=\scriptsize, align=center]{TCP\\JSON} (gu.west);
	\end{tikzpicture}
	\caption{Sơ đồ gói các nhóm mô-đun hai phía và quan hệ phụ thuộc (bản nháp)}
	\label{fig:packages}
\end{figure}

\section{Thiết kế chi tiết}

\subsection{Máy trạng thái điều khiển phương tiện}

Điểm thiết kế cốt lõi của phần tương tác là một máy trạng thái xác định quyền
điều khiển và hành vi của mỗi phương tiện. Mỗi phương tiện luôn ở một trong bốn
trạng thái: được \gls{SUMO} điều khiển, được người dùng điều khiển, trở thành xác
xe sau va chạm, hoặc bị hủy. Quyền điều khiển và việc bật/tắt vật lý được quyết
định hoàn toàn bởi trạng thái này, thay vì dựa vào sự hiện diện của thành phần.
Hình \ref{fig:existstate} mô tả máy trạng thái.

\begin{figure}[H]
	\centering
	\begin{tikzpicture}[every node/.style={font=\small}, >=Latex,
		st/.style={draw, rounded corners, align=center, minimum width=2.6cm, minimum height=0.9cm}]
		\node[st] (server) {Do SUMO\\điều khiển};
		\node[st, right=3.5cm of server] (client) {Do người dùng\\điều khiển};
		\node[st, below=2.2cm of client] (wreck) {Xác xe};
		\node[st, below=2.2cm of server] (dead) {Bị hủy};
		\draw[->] (server) -- node[above, font=\scriptsize, align=center] {chiếm\\quyền} (client);
		\draw[->] (client.north) to[out=120,in=60,looseness=1.4] node[above, font=\scriptsize, align=center] {trả quyền (đưa về làn)} (server.north);
		\draw[->] (server) -- node[left, font=\scriptsize, align=center] {bị tông} (wreck);
		\draw[->] (client) -- node[right, font=\scriptsize, align=center] {trả quyền\\khi đã va chạm} (wreck);
		\draw[->] (wreck) -- node[below, font=\scriptsize, align=center] {sau N bước} (dead);
		\draw[->] (client.south west) -- node[above, font=\scriptsize, align=center, sloped] {thả quá xa làn} (dead.north east);
		\draw[->] (dead.south) to[out=-90,in=-90,looseness=1.2] node[below, font=\scriptsize] {tái sử dụng} (server.south);
	\end{tikzpicture}
	\caption{Máy trạng thái điều khiển phương tiện (bản nháp)}
	\label{fig:existstate}
\end{figure}

Khi một phương tiện do \gls{SUMO} điều khiển, vị trí của nó được cập nhật theo dữ
liệu mô phỏng và phần vật lý được tắt. Khi người dùng chiếm quyền, hệ thống bật
vật lý để người dùng lái; vị trí xe được phản ánh ngược về \gls{SUMO} để các xe
khác phản ứng. Khi một xe đang do người dùng điều khiển hoặc một xác xe va chạm
với xe do \gls{SUMO} điều khiển, xe bị tông chuyển thành xác xe, bị đông cứng
trong mô phỏng để gây ùn ứ, và tự bị hủy sau một số bước định trước. Khi trả
quyền mà xe không va chạm, hệ thống đưa xe trở lại làn gần nhất; nếu quá xa, xe bị
hủy. Trạng thái bị hủy đưa phương tiện về bộ nhớ tái sử dụng. Chi tiết cơ chế
chuyển đổi vật lý và đưa xe về làn được trình bày trong Chương 5.

\subsection{Hai luồng dữ liệu giữa mô phỏng và hiển thị}

Hệ thống có hai luồng dữ liệu ngược chiều. Luồng tải xuống truyền trạng thái mô
phỏng từ Python sang Unity theo từng bước, gồm chỉ số bước hiện tại, trạng thái
đèn tín hiệu và danh sách đối tượng giao thông; trong đó mỗi đối tượng chứa định
danh, loại, vị trí, hướng và một số cờ trạng thái. Luồng tải lên truyền dữ liệu
các phương tiện do người dùng điều khiển từ Unity về Python theo lô trong một
thông điệp, gồm định danh, vị trí, hướng, tốc độ và trạng thái tồn tại của xe.
Phía Python dùng luồng tải lên để phản ánh xe người lái vào \gls{SUMO} và để đối
soát danh sách phương tiện đang được quản lý. Trình tự trao đổi thông điệp trong
mỗi bước mô phỏng được minh họa ở Hình \ref{fig:dataflow-seq}.

\begin{figure}[H]
	\centering
	\begin{tikzpicture}[every node/.style={font=\footnotesize}, >=Latex]
		% header boxes
		\node[draw, rounded corners, minimum width=3.2cm, minimum height=0.8cm, align=center] (py) {Python\\(mô phỏng)};
		\node[draw, rounded corners, minimum width=3.2cm, minimum height=0.8cm, align=center, right=7cm of py] (un) {Unity\\(hiển thị \& tương tác)};
		% lifelines
		\draw[dashed] (py.south) -- ++(0,-5.4) coordinate (pyEnd);
		\draw[dashed] (un.south) -- ++(0,-5.4) coordinate (unEnd);
		% loop frame
		\draw[rounded corners, draw=black!55] ($(py.south)+(-1.5,-0.5)$) rectangle ($(un.south)+(1.5,-5.1)$);
		\node[anchor=north west, font=\scriptsize\bfseries, fill=black!8, inner sep=2pt]
			at ($(py.south)+(-1.5,-0.5)$) {lặp [mỗi bước mô phỏng]};
		% messages
		\draw[-{Latex}] ($(py.south)+(0,-1.3)$) -- node[above, font=\scriptsize, align=center]
			{5050: trạng thái bước\\(đèn tín hiệu, danh sách đối tượng)} ($(un.south)+(0,-1.3)$);
		\draw[{Latex}-] ($(py.south)+(0,-3.0)$) -- node[above, font=\scriptsize, align=center]
			{5053: lô dữ liệu xe người lái\\(vị trí, hướng, tốc độ)} ($(un.south)+(0,-3.0)$);
		\draw[{Latex}-] ($(py.south)+(0,-4.4)$) -- node[above, font=\scriptsize]
			{5054: lệnh điều khiển (tạm dừng / tiếp tục / kết thúc)} ($(un.south)+(0,-4.4)$);
	\end{tikzpicture}
	\caption{Trình tự trao đổi thông điệp giữa mô phỏng và hiển thị trong một bước (bản nháp)}
	\label{fig:dataflow-seq}
\end{figure}

\subsection{Thiết kế một số lớp chủ đạo}

Trong phạm vi báo cáo, một số lớp chủ đạo được trình bày như sau.

\begin{itemize}[leftmargin=2em]
	\item \texttt{TrafficerManager} (Unity): quản lý vòng đời phương tiện theo dữ
	liệu tải xuống, tạo, cập nhật và thu hồi đối tượng; đếm tập trung số bước để
	hủy xác xe nhằm tránh lỗi khi đối tượng bị tạm ẩn.
	\item \texttt{UnityVehicle} (Unity): xử lý chiếm quyền, trả quyền và chuyển xe
	thành xác xe; bật/tắt vật lý theo trạng thái; thu thập dữ liệu xe để gửi lên.
	\item \texttt{WheelController} (Unity): điều khiển bốn bánh xe gồm tăng tốc,
	đánh lái và phanh khi xe ở chế độ lái.
	\item \texttt{VehicleSender} (Unity): gửi dữ liệu phương tiện theo lô, giữ lô
	mới nhất để truyền lên server.
	\item \texttt{SimulationSession} (Python): gom toàn bộ kết quả của một lần chạy
	vào một thư mục phiên và ghi chuỗi trạng thái theo bước dưới dạng luồng.
\end{itemize}

Quan hệ giữa các lớp chủ đạo phía Unity được minh họa ở Hình \ref{fig:class-core}.

\begin{figure}[H]
	\centering
	\begin{tikzpicture}[
		every node/.style={font=\footnotesize}, >=Latex,
		class/.style={rectangle split, rectangle split parts=2, draw, rounded corners=2pt,
			align=left, text width=3.5cm, inner sep=4pt}]
		\node[class] (tm) {\textbf{TrafficerManager}
			\nodepart{second} quản lý vòng đời xe\\ tạo / cập nhật / thu hồi\\ đếm bước hủy xác xe};
		\node[class, right=3.2cm of tm] (uv) {\textbf{UnityVehicle}
			\nodepart{second} chiếm / trả quyền\\ bật / tắt vật lý\\ thu thập dữ liệu xe};
		\node[class, below=2.1cm of uv] (wc) {\textbf{WheelController}
			\nodepart{second} tăng tốc\\ đánh lái, phanh};
		\node[class, below=2.1cm of tm] (vs) {\textbf{VehicleSender}
			\nodepart{second} gửi dữ liệu xe theo lô\\ giữ lô mới nhất};
		% quan he
		\draw[{Diamond[open]}-{Latex}] (tm.east) -- node[above, font=\scriptsize]{quản lý} node[below, font=\scriptsize]{1..*} (uv.west);
		\draw[{Diamond}-] (uv.south) -- node[right, font=\scriptsize]{sở hữu} (wc.north);
		\draw[-{Latex}] (uv.south west) -- node[above, font=\scriptsize, sloped]{gửi qua} (vs.north east);
	\end{tikzpicture}
	\caption{Sơ đồ lớp các lớp chủ đạo phía Unity (bản nháp). Lớp \texttt{SimulationSession} thuộc phía Python, không liên kết trực tiếp với nhóm lớp này.}
	\label{fig:class-core}
\end{figure}

\subsection{Lưu trữ dữ liệu}

Hệ thống không dùng cơ sở dữ liệu truyền thống mà lưu dữ liệu dưới dạng tệp. Dữ
liệu cấu hình và mạng mô phỏng là các tệp của \gls{SUMO}. Dữ liệu kết quả của mỗi
lần chạy được gom vào một thư mục phiên đặt theo tên bản đồ và thời điểm chạy, gồm
nhật ký hành trình của xe, bản tóm tắt, dữ liệu mạng đường và chuỗi trạng thái
theo từng bước. Cách lưu này phù hợp với mục tiêu mô phỏng, hiển thị và thu thập
thống kê theo từng lần chạy, đồng thời cho phép phát lại. Cấu trúc một thư mục
phiên tiêu biểu được minh họa ở Hình \ref{fig:session-tree}.

\begin{figure}[H]
	\centering
	\begin{tikzpicture}[>=Latex,
		fname/.style={font=\footnotesize\ttfamily, anchor=west, inner sep=1pt},
		desc/.style={font=\footnotesize, anchor=west, inner sep=1pt, text=black!60}]
		% nodes
		\node[fname] (root) at (0,0) {result/};
		\node[fname] (sess) at (0.9,-0.75) {HoanKiem-2026-06-13-09-30-05/};
		\node[fname] (c1) at (1.8,-1.5) {trips.csv};
		\node[fname] (c2) at (1.8,-2.2) {summary.json};
		\node[fname] (c3) at (1.8,-2.9) {road\_data.json};
		\node[fname] (c4) at (1.8,-3.6) {scenario.json};
		% descriptions
		\node[desc] at (5.6,-1.5) {nhật ký hành trình của xe};
		\node[desc] at (5.6,-2.2) {bản tóm tắt lần chạy};
		\node[desc] at (5.6,-2.9) {dữ liệu mạng đường (nút, cạnh, vạch qua đường)};
		\node[desc] at (5.6,-3.6) {chuỗi trạng thái mô phỏng theo từng bước};
		% connectors
		\draw (0.35,-0.25) |- (sess.west);
		\draw (1.25,-1.0) |- (c1.west);
		\draw (1.25,-1.0) |- (c2.west);
		\draw (1.25,-1.0) |- (c3.west);
		\draw (1.25,-1.0) |- (c4.west);
	\end{tikzpicture}
	\caption{Cấu trúc một thư mục phiên tiêu biểu (bản nháp), đặt tên theo bản đồ và thời điểm chạy. Tệp \texttt{trips.csv} đổi thành \texttt{trips-ped.csv} khi kịch bản có người đi bộ.}
	\label{fig:session-tree}
\end{figure}

\subsection{Thiết kế giao diện người dùng}
\label{subsec:thiet-ke-giao-dien}

Khác với các ứng dụng hướng biểu mẫu, giao diện của hệ thống được thiết kế xoay
quanh việc \emph{quan sát và tương tác với cảnh ba chiều}, nên thành phần trung
tâm không phải là các ô nhập liệu mà là khung nhìn (camera) cùng một lớp điều
khiển gọn nhẹ phủ lên trên. Giao diện được tổ chức theo nguyên tắc: ưu tiên
không gian hiển thị cho cảnh, các bảng điều khiển chỉ hiện khi cần và có thể ẩn
đi; thao tác chính dùng chuột và bàn phím; và các nút mang tính ngữ cảnh, chỉ
xuất hiện khi thực sự áp dụng được cho đối tượng đang chọn.

\textbf{Chế độ hiển thị.} Hệ thống có hai chế độ với bố cục giao diện riêng: chế
độ thời gian thực (hiển thị mô phỏng đang chạy, kèm các lệnh điều khiển quá
trình chạy) và chế độ phát lại (đọc lại một phiên đã lưu). Trước khi vào cảnh,
người dùng chọn chế độ, bản đồ và tham số mô phỏng ở bước cấu hình khởi chạy.

\textbf{Camera và góc nhìn.} Camera hỗ trợ hai góc nhìn nhằm phục vụ cả nhu cầu
giám sát tổng thể lẫn trải nghiệm của người tham gia giao thông:
\begin{itemize}[leftmargin=2em]
	\item \textbf{Góc nhìn tự do}: camera bay tự do trên cao, di chuyển bằng phím
	(\texttt{WASD}) và xoay bằng chuột, với tốc độ di chuyển và độ nhạy chuột
	điều chỉnh được. Đây là góc nhìn để quan sát toàn cảnh mạng đường và dòng
	giao thông.
	\item \textbf{Góc nhìn gắn vào phương tiện}: camera gắn vào một phương tiện
	cụ thể để nhìn theo nó. Người dùng chọn phương tiện bằng cách trỏ và bấm
	(point-and-select), hoặc chọn ngẫu nhiên một phương tiện đang được hiển thị trong danh sách. 
	Đây là góc nhìn tạo nên trải nghiệm theo người tham
	gia giao thông và là tiền đề cho thao tác chiếm quyền điều khiển.
\end{itemize}
Trạng thái camera hiện hành (tự do hay đang gắn vào phương tiện nào) được hiển
thị thường trực để người dùng định hướng.

\textbf{Lớp điều khiển và bảng thông tin (HUD).} Phủ trên khung nhìn là một lớp
giao diện gọn nhẹ. Góc dưới bên trái là bảng thông tin hiển thị trạng thái
camera (tự do hay đang gắn vào phương tiện nào) cùng các chỉ số tổng quan của
phiên chạy: độ trễ đầu--cuối, số phương tiện và số người đi bộ. Góc dưới bên
phải là hàng nút điều khiển chính gồm: nút \emph{chiếm/trả quyền điều khiển}
mang tính ngữ cảnh (chỉ hiện khi camera đang gắn vào một phương tiện có thể lái,
và tự đổi nhãn giữa ``Chiếm quyền'' và ``Trả quyền'' theo trạng thái xe); nút
\emph{tạm dừng/tiếp tục} mô phỏng; nút \emph{Tạo xe Unity} để sinh một phương
tiện riêng do người dùng điều khiển (và đổi thành ``Hủy xe Unity'' để loại bỏ
nó); và nút \emph{Người/Xe} để ẩn/hiện danh sách các đối tượng giao thông
(phương tiện và người đi bộ) đang có trong cảnh theo mã định danh, chọn một mục
trong danh sách sẽ gắn camera vào đúng đối tượng tương ứng. Hình
\ref{fig:ui-layout} minh họa bố cục giao diện trong cảnh.

Ngoài hàng nút chính, các thiết lập ít dùng hơn được đặt trong các bảng tùy
chọn mở bằng phím \texttt{Esc}, gồm: điều chỉnh tốc độ mô phỏng, tốc độ và độ
nhạy của camera, tùy chọn lọc theo vùng quan sát phục vụ tối ưu hiển thị, và
lệnh kết thúc phiên chạy. Khi một bảng tùy chọn mở, con trỏ chuột được mở khóa
để thao tác và được khóa lại khi đóng để tiếp tục điều khiển camera; ngoài ra
bảng hướng dẫn phím điều khiển cũng có thể bật/tắt khi cần.

\begin{figure}[H]
	\centering
	\resizebox{0.98\linewidth}{!}{%
	\begin{tikzpicture}[>=Latex, font=\footnotesize,
		btn/.style={draw=black!45, rounded corners=3pt, fill=black!12,
			minimum height=0.6cm, inner xsep=6pt, align=center, font=\scriptsize},
		item/.style={draw=black!40, fill=white, rounded corners=1pt,
			minimum width=1.95cm, minimum height=0.44cm, inner sep=1pt,
			font=\scriptsize\ttfamily},
		anno/.style={font=\scriptsize\itshape, text=black!55}]
		\def\W{16.5} \def\Hh{9.3}
		% --- Khung nhin ba chieu (viewport): nen dong nhat, khong co duong chan troi ---
		\fill[blue!8] (0,0) rectangle (\W,\Hh);
		\draw[black!55] (0,0) rectangle (\W,\Hh);              % viền khung nhìn
		% --- Tam ngam (crosshair) o giua khung nhin ---
		\coordinate (c) at (\W/2,\Hh/2);
		\draw[black!70, line width=0.6pt] ($(c)+(-0.30,0)$) -- ($(c)+(-0.10,0)$);
		\draw[black!70, line width=0.6pt] ($(c)+(0.10,0)$)  -- ($(c)+(0.30,0)$);
		\draw[black!70, line width=0.6pt] ($(c)+(0,-0.30)$) -- ($(c)+(0,-0.10)$);
		\draw[black!70, line width=0.6pt] ($(c)+(0,0.10)$)  -- ($(c)+(0,0.30)$);
		% --- Bang thong tin & trang thai camera (goc duoi trai, tu co vua khit) ---
		\node[anchor=south west, text=white, fill=black!58, rounded corners=2pt,
			inner sep=5pt, align=left, font=\scriptsize] (info) at (0.15,0.2)
			{{\footnotesize\bfseries Trạng thái camera: Tự do}\\[2pt]
			 Độ trễ: 10\,ms \quad Số xe: 100 \quad Số người: 1000};
		% --- Lop dieu khien (goc duoi phai): neo tu mep phai, noi lien co khe ho ---
		\node[btn, anchor=east] (b4) at (\W-0.2,0.85) {Người/Xe};
		\node[btn, anchor=east] (b3) at ([xshift=-0.18cm]b4.west) {Tạo xe Unity};
		\node[btn, anchor=east] (b2) at ([xshift=-0.18cm]b3.west) {Tạm dừng};
		\node[btn, anchor=east] (b1) at ([xshift=-0.18cm]b2.west) {Chiếm quyền};
		% --- Danh sach Trafficer (hien khi bam Nguoi/Xe), xep ngay tren cac nut ---
		\foreach \name [count=\i from 0] in {f\_1.0, f\_10.0, f\_14.0, f\_15.0, f\_16.0, f\_2.0} {
			\node[item, anchor=south east] at (\W-0.2, 1.35 + \i*0.5) {\name};
		}
		% --- Chu thich ---
		\node[anno, anchor=west] (a1) at (\W/2+0.55,\Hh/2+0.5)
			{tâm ngắm: trỏ--chọn phương tiện};
		\draw[black!45, ->] (a1.west) -- ($(c)+(0.34,0.16)$);
		\node[anno, anchor=south west] at (0.15,1.95) {bảng thông tin \& trạng thái camera};
		\node[anno, anchor=south west] at ($(b1.north west)+(0,0.12)$) {lớp điều khiển};
		% chu thich danh sach Trafficer
		\node[anno, anchor=east, align=right] (la) at (13.8,3.1)
			{\emph{Người/Xe}: mở/đóng danh sách Trafficer;\\chọn một mục để gắn camera vào đối tượng};
		\draw[black!45, ->] (la.east) -- (\W-2.25,3.1);
	\end{tikzpicture}}
	\caption{Bố cục giao diện người dùng trong cảnh ba chiều (bản nháp)}
	\label{fig:ui-layout}
\end{figure}

\section{Xây dựng ứng dụng}

\subsection{Thư viện và công cụ sử dụng}

Bảng \ref{tab:tools} liệt kê các công cụ chính được sử dụng trong dự án.

\begin{table}[H]
	\centering
	\caption{Danh sách thư viện và công cụ sử dụng}
	\label{tab:tools}
	\begin{tabular}{p{0.34\textwidth} p{0.30\textwidth} p{0.26\textwidth}}
		\hline
		\textbf{Mục đích} & \textbf{Công cụ} & \textbf{Phiên bản} \\
		\hline
		Mô phỏng giao thông vi mô & Eclipse SUMO & 1.24.0 \\
		Điều khiển mô phỏng & TraCI (Python) & đi kèm SUMO \\
		Ngôn ngữ điều phối & Python & 3.x \\
		Hiển thị 3D và tương tác & Unity Editor & 6000.3.9f1 \\
		Phân tích JSON trong Unity & Newtonsoft.Json & 3.2.x \\
		Môi trường phát triển & Visual Studio Code & --- \\
		Quản lý phiên bản & Git & --- \\
		Hệ điều hành thử nghiệm & Windows & 10/11 \\
		\hline
	\end{tabular}
\end{table}

\subsection{Kết quả đạt được và thống kê mã nguồn}

Sản phẩm chính gồm: bộ sinh mạng và tuyến cho \gls{SUMO} (từ lưới mê cung hoặc từ
bản đồ OpenStreetMap); ứng dụng Unity hiển thị mô phỏng ba chiều có tương tác lái
xe; và bộ lưu trữ phiên cùng nhật ký kết quả. Bảng \ref{tab:loc} thống kê nhanh
quy mô mã nguồn do nhóm tự viết, không tính thư viện bên thứ ba.

\begin{table}[H]
	\centering
	\caption{Thống kê quy mô mã nguồn}
	\label{tab:loc}
	\begin{tabular}{p{0.40\textwidth} c c}
		\hline
		\textbf{Thành phần} & \textbf{Số tệp} & \textbf{Số dòng} \\
		\hline
		Python (mô phỏng và điều phối) & 40 & 5775 \\
		Unity (mã nguồn C\#) & 100 & 8060 \\
		\hline
	\end{tabular}
\end{table}

\subsection{Minh họa các chức năng chính}

Các chức năng chính đã hiện thực gồm: sinh mạng đường từ lưới mê cung hoặc bản đồ
thực tế; chạy mô phỏng và đồng bộ trạng thái sang Unity theo thời gian thực; hiển
thị phương tiện, đèn tín hiệu và người đi bộ trong không gian ba chiều; chiếm
quyền điều khiển và lái xe thủ công; xử lý va chạm sinh xác xe và trả quyền đưa xe
về làn; cùng các thao tác điều khiển quá trình chạy. Các hình minh họa dưới đây
thể hiện một số chức năng tiêu biểu.

\begin{figure}[H]
	\centering
	\includegraphics[width=0.92\linewidth]{Hinhve/hinh4.8-anh-goc-nhin-3d-tu-tren-cao.png}
	\caption{Minh họa cảnh giao thông ba chiều trong Unity: góc nhìn từ trên cao}
	\label{fig:demo-scene}
\end{figure}

\begin{figure}[H]
	\centering
	\begin{subfigure}[b]{0.85\linewidth}
		\includegraphics[width=\linewidth]{Hinhve/hinh4.9.2-goc-nhin-cua-1-xe-tu-hanh-boi-may-chu.png}
		\caption{Quan sát xe tự hành, nút ``Chiếm quyền'' hiện ra}
		\label{fig:demo-drive-observe}
	\end{subfigure}
	\vspace{0.5em}
	\begin{subfigure}[b]{0.85\linewidth}
		\includegraphics[width=\linewidth]{Hinhve/hinh4.9-goc-nhin-trong-xe.png}
		\caption{Đang lái xe (camera gắn CLIENT\_CAR)}
		\label{fig:demo-drive-active}
	\end{subfigure}
	\caption{Minh họa chế độ lái xe thủ công: từ quan sát đến chiếm quyền điều khiển}
	\label{fig:demo-drive}
\end{figure}

\begin{figure}[H]
	\centering
	\includegraphics[width=0.92\linewidth]{Hinhve/hinh4.10-va-cham-gay-un-tac.png}
	\caption{Minh họa xử lý va chạm: xác xe nằm lại gây ùn ứ dòng giao thông}
	\label{fig:demo-crash}
\end{figure}

\section{Kiểm thử}

Do hệ thống có tính tích hợp cao giữa \gls{SUMO} và Unity, việc kiểm thử được
thực hiện theo hai hướng: kiểm thử chức năng thủ công theo use case và kiểm thử
hộp đen cho các mô-đun xử lý dữ liệu và giao thức. Bảng \ref{tab:testcases} liệt
kê một số trường hợp kiểm thử tiêu biểu.

\begin{table}[H]
	\centering
	\caption{Một số trường hợp kiểm thử tiêu biểu}
	\label{tab:testcases}
	\begin{tabular}{p{0.08\textwidth} p{0.34\textwidth} p{0.46\textwidth}}
		\hline
		\textbf{Mã} & \textbf{Mục tiêu} & \textbf{Dữ liệu vào / Kỳ vọng} \\
		\hline
		TC01 & Tách khung thông điệp TCP & Gửi hai thông điệp liên tiếp trong một lần gửi; phía nhận tách đúng hai thông điệp. \\
		TC02 & Nhận dữ liệu mạng đường một lần & Unity yêu cầu dữ liệu mạng; nhận đủ và dựng cảnh đúng. \\
		TC03 & Hiển thị trạng thái đèn & Mô phỏng trả về chuỗi trạng thái đèn; Unity hiển thị đúng đỏ/vàng/xanh. \\
		TC04 & Chiếm quyền và lái xe & Người dùng chiếm một xe; xe chuyển sang lái bằng vật lý, các xe khác phản ứng theo. \\
		TC05 & Va chạm sinh xác xe & Xe người lái tông xe của mô phỏng; xe bị tông thành xác xe và biến mất sau N bước. \\
		TC06 & Trả quyền và đưa về làn & Thả xe gần làn: xe chạy lại tự động; thả xe quá xa làn: xe bị hủy. \\
		\hline
	\end{tabular}
\end{table}

Trong điều kiện chạy thử nghiệm, các trường hợp trên đạt yêu cầu. Hiện tượng lệch
đồng bộ chủ yếu xuất hiện khi tốc độ mô phỏng đặt quá cao; khi điều chỉnh tốc độ
hợp lý, quan sát ổn định hơn.

\section{Thực nghiệm và đánh giá}

\subsection{Môi trường và kịch bản thực nghiệm}

Hệ thống được thử nghiệm trên máy cá nhân chạy Windows với hai tiến trình: tiến
trình Python kết hợp \gls{SUMO} và tiến trình Unity. Hai nhóm kịch bản được sử
dụng: nhóm kịch bản sinh từ lưới mê cung với kích thước tăng dần nhằm khảo sát
khả năng chịu tải theo số lượng đối tượng; và nhóm kịch bản từ bản đồ thực tế
OpenStreetMap nhằm kiểm chứng khả năng dựng cảnh và hiển thị trên dữ liệu thật.

\subsection{Phương pháp đánh giá}

Việc đánh giá tập trung vào ba khía cạnh. Thứ nhất là tính đúng đắn của hiển thị,
được kiểm chứng bằng cách đối chiếu cảnh ba chiều trong Unity với giao diện hai
chiều của \gls{SUMO} trên cùng một kịch bản. Thứ hai là hiệu năng hiển thị, đo
bằng tốc độ khung hình khi số lượng phương tiện thay đổi. Thứ ba là tính ổn định
của tương tác, quan sát qua các thao tác chiếm quyền, lái xe, va chạm và trả
quyền.

\subsection{Kết quả}

% [GHI CHÚ] Cần điền số liệu đo thực tế (tốc độ khung hình theo số lượng xe, độ
% trễ đồng bộ) và ảnh chụp đối chiếu. Bảng dưới là khung để điền kết quả.
Bảng \ref{tab:perf} là khung ghi nhận kết quả hiệu năng theo quy mô kịch bản; các
giá trị cụ thể được điền sau khi đo trên môi trường thực nghiệm.

\begin{table}[H]
	\centering
	\caption{Khung ghi nhận hiệu năng theo quy mô kịch bản}
	\label{tab:perf}
	\begin{tabular}{p{0.30\textwidth} c c c}
		\hline
		\textbf{Kịch bản} & \textbf{Số phương tiện} & \textbf{Tốc độ khung hình} & \textbf{Ghi chú} \\
		\hline
		Mê cung nhỏ & --- & --- & --- \\
		Mê cung trung bình & --- & --- & --- \\
		Mê cung lớn & --- & --- & --- \\
		Bản đồ OpenStreetMap & --- & --- & --- \\
		\hline
	\end{tabular}
\end{table}

Trên cơ sở quan sát trong quá trình phát triển, hệ thống dựng được mạng đường từ
cả hai nguồn dữ liệu, tái hiện chuyển động của phương tiện và người đi bộ một cách
ổn định, và cho phép người dùng tương tác lái xe trong cảnh. Việc lọc đối tượng
theo vùng quan sát và tái sử dụng đối tượng giúp duy trì hiệu năng khi số lượng
phương tiện tăng. Các số liệu định lượng chi tiết sẽ được bổ sung sau khi hoàn
tất quá trình đo trên môi trường thực nghiệm.

\subsection{Triển khai}

Trong phạm vi đồ án, hệ thống chạy theo mô hình một máy và một người dùng. Quy
trình chạy gồm các bước: chuẩn bị kịch bản và khởi chạy tiến trình Python kết hợp
\gls{SUMO} thông qua giao diện khởi chạy; sau đó chạy ứng dụng Unity để kết nối,
nhận dữ liệu và hiển thị. Ba kênh TCP được dùng tương ứng cho dữ liệu mô phỏng, dữ liệu
phương tiện do người dùng lái và lệnh điều khiển. Kết quả mỗi lần chạy được lưu
cục bộ theo từng phiên để phục vụ phân tích và phát lại.

Chương này đã trình bày thiết kế kiến trúc ba tầng với ba kênh TCP và hai luồng
dữ liệu, máy trạng thái điều khiển phương tiện, thiết kế các mô-đun và lớp chủ
đạo, cùng phương pháp kiểm thử, thực nghiệm và triển khai. Các giải pháp kỹ thuật
nổi bật phát sinh trong quá trình thiết kế được trình bày chi tiết trong Chương 5.

\end{document}
